\section{线性数据结构章正文案例推导补强}

这一节补齐本章正文出现过的 LeetCode 案例推导。阅读顺序统一按“题目说明、样例入口、暴力基线、暴力过程、重复工作、优化条件、相邻状态、状态复用、指针和字段、代码落点、正确性、边界实验、复杂度变化、迁移追问”展开；目的不是新增题量，而是把正文中的案例都讲成可复述、可证明、可迁移的解题链。

\subsection{LeetCode 438 找到字符串中所有字母异位词}

\textbf{题目说明。} LeetCode 438 找到字符串中所有字母异位词 的题面入口要先还原成具体任务：固定长度窗口内字符频次等于目标频次即为答案。

\textbf{样例入口。} 拿 s=cbaebabacd、p=abc 手算，长度为 3 的窗口 cba 频次和 abc 相同，起点 0 是答案；窗口右移时只删除 c、加入 e。

\textbf{暴力基线。} 先统计 p 的频次 target，再枚举 s 中每个长度 k=p.size() 的窗口，为每个窗口重新统计 window 后比较。

\textbf{暴力过程。} 以 s=cbaebabacd、p=abc 为例，暴力窗口依次是 cba、bae、aeb、eba、bab、aba、bac、acd。每个窗口都重新数 k 个字符。

\textbf{重复工作。} 题面模型：固定长度窗口内字符频次等于目标频次即为答案。暴力版的慢点要落到本题：相邻窗口只差一个出去字符和一个进来字符。从 cba 到 bae，只有 c 出去、e 进来，b 和 a 不该重数。后面的优化只消掉这类重复工作，不能改变答案集合。

\textbf{优化条件。} 本题的关键观察是：右进左出维护计数，窗口长度达到目标长度后检查差异。这句话必须能翻译成可维护状态；如果题目条件改变后观察不再成立，就不能继续套原来的写法。

\textbf{相邻状态。} 规律是固定长度窗口维护字符频次差，重叠答案不能漏。把这个特例继续拆开看：右端进入只带来一个新元素，左端离开只删掉一个旧元素；只有当旧左端被移走后不可能再产生更优答案时，窗口才可以单向推进。若这个单调淘汰条件不存在，就要换成前缀和、队列或其他状态。

\textbf{状态复用。} k 来自 p.length()。先统计第一个窗口，之后每次窗口右移一格：window[s[i-1]]--，window[s[i+k-1]]++。手算时只改被新输入影响的那一小部分，而不是回头重算完整候选。

\textbf{指针和字段。} 状态字段要能说清：i 是当前窗口起点，out=s[i-1] 是离开字符，in=s[i+k-1] 是进入字符，target/window 是 26 长度频次数组。手算时按这个协议跟踪：第一个窗口 cba 与 abc 频次相同，记录 0；滑到 bae 后 c 计数减一、e 加一；滑到 bac 时频次再次相同，记录 6。

\textbf{代码落点。} 关键代码是固定窗口右进左出。若用 unordered\_map，计数减到 0 的 key 要 erase；若题目限定小写字母，用 array/vector 长度 26 更简单。关键代码行必须能逐句翻译成“旧状态怎样变成新状态”，否则就只是模板。

\textbf{正确性。} 每个被检查窗口长度都等于 k，window 频次始终精确表示当前窗口；频次等于 target 当且仅当该窗口是 p 的异位词。

\textbf{边界实验。} 先用目标比源串长、重复字符、字符集固定、答案重叠做反例检查。实验时重点跑目标比源串长、重复字符、字符集固定、答案重叠，并打印左右端、窗口计数和答案。若左端发生回退，或者窗口失去合法性后仍更新答案，就能很快定位错误。

\textbf{复杂度变化。} 暴力 O(nk)，固定窗口维护 O(n*26)，因字符集固定通常记作 O(n)，空间 O(26)。

\textbf{迁移追问。} 如果题目改成“维护差异计数可以避免每次比较整个数组”，先判断原状态是否仍然覆盖新答案集合，再决定是微调边界、改 value 语义，还是换算法。

\subsection{LeetCode 206 反转链表}

\textbf{题目说明。} LeetCode 206 反转链表 的题面入口要先还原成具体任务：每次把当前节点的 next 指向已经反转好的前缀头。

\textbf{样例入口。} 拿 1->2->3 手算，处理 1 时先保存 next=2，再让 1->null；处理 2 时保存 3，再让 2->1。

\textbf{暴力基线。} 慢版本可以先把节点顺序抄到数组里，按数组推导目标顺序。回到链表后，难点不再是值的顺序，而是节点身份和 next 指针不能丢。放到本题就是：先按“每次把当前节点的 next 指向已经反转好的前缀头”完整试慢版本；样例里已经能看到“规律是 prev 保存已反转前缀头，cur 保存待处理节点，每次改 next 前必须保存后继”，所以优化前必须先能说清慢版本枚举了哪些候选。

\textbf{暴力过程。} 拿 1->2->3 手算，处理 1 时先保存 next=2，再让 1->null；处理 2 时保存 3，再让 2->1。

\textbf{重复工作。} 题面模型：每次把当前节点的 next 指向已经反转好的前缀头。暴力版的慢点要落到本题：慢版本会围绕“每次把当前节点的 next 指向已经反转好的前缀头”借助数组或多次扫描确认目标顺序。样例规律是“规律是 prev 保存已反转前缀头，cur 保存待处理节点，每次改 next 前必须保存后继”，说明优化点不在节点值，而在少量指针角色能否保持已处理链、未处理链和接回位置都不丢。后面的优化只消掉这类重复工作，不能改变答案集合。

\textbf{优化条件。} 本题的关键观察是：先保存后继，再改指针，再推进 previous 和 current。这句话必须能翻译成可维护状态；如果题目条件改变后观察不再成立，就不能继续套原来的写法。

\textbf{相邻状态。} 规律是 prev 保存已反转前缀头，cur 保存待处理节点，每次改 next 前必须保存后继。把这个特例继续拆开看：每个指针都要有角色，上一段尾、当前段头、当前节点和后继入口不能混。只要一次改 next 之前没有保存后继，就可能丢掉整段未处理链。

\textbf{状态复用。} 把数组辅助结构换成几个稳定指针角色；重连顺序必须保证已处理链合法、未处理链仍可达。本题落点是：规律是 prev 保存已反转前缀头，cur 保存待处理节点，每次改 next 前必须保存后继。手算时只改被新输入影响的那一小部分，而不是回头重算完整候选。

\textbf{指针和字段。} 状态字段要能说清：虚拟头、上一段尾、当前段头、当前节点、后继节点；任何改 next 的操作之前都要保存后续入口。手算时按这个协议跟踪：拿 1->2->3 手算，处理 1 时先保存 next=2，再让 1->null；处理 2 时保存 3，再让 2->1。然后按协议复查：手算时画出 prev、cur、next 和下一段头节点。每次改 next 前先保存后继，这是链表推导的安全线。

\textbf{代码落点。} 所有重连步骤按保存后继、断开或反转、接回前缀、接回后缀的顺序写；最后检查是否形成环或丢节点。关键代码行必须能逐句翻译成“旧状态怎样变成新状态”，否则就只是模板。

\textbf{正确性。} 证明从链结构开始：已处理链连通、未处理链可达、二者连接点明确，节点没有丢失也没有成环。

\textbf{边界实验。} 先用空链、单节点、两个节点、不要丢失剩余链做反例检查。实验时覆盖空链、单节点、两个节点、不要丢失剩余链，每步画出前驱、当前、后继和下一段头。链表题不要只看输出值，还要检查节点是否丢失或成环。

\textbf{复杂度变化。} 暴力可能多次扫描链表或拷贝到数组；优化后每个节点通常只被访问、重连常数次，目标是线性时间和常数额外空间。

\textbf{迁移追问。} 如果题目改成“K 组反转是在这段局部反转外增加分组边界”，先判断原状态是否仍然覆盖新答案集合，再决定是微调边界、改 value 语义，还是换算法。

\subsection{LeetCode 21 合并两个有序链表}

\textbf{题目说明。} LeetCode 21 合并两个有序链表 的题面入口要先还原成具体任务：两个链表已经各自有序，下一节点只能来自两个当前头中较小者。

\textbf{样例入口。} 拿 1->3 和 2->4 手算，先接 1，再比较 3 和 2 接 2，之后接 3 和 4。

\textbf{暴力基线。} 暴力可以把两个链表的值全部放入数组，排序后重新建链表。

\textbf{暴力过程。} 以 1->3 和 2->4 为例，数组法得到 [1, 2, 3, 4]；但原链表本来有序，每一步只需要比较两个当前头。

\textbf{重复工作。} 题面模型：两个链表已经各自有序，下一节点只能来自两个当前头中较小者。暴力版的慢点要落到本题：浪费在完整排序和新建节点。两个链表各自有序，下一节点一定来自 l1 或 l2 的较小头节点。后面的优化只消掉这类重复工作，不能改变答案集合。

\textbf{优化条件。} 本题的关键观察是：虚拟头统一结果链表，尾指针不断接上较小节点并推进来源链。这句话必须能翻译成可维护状态；如果题目条件改变后观察不再成立，就不能继续套原来的写法。

\textbf{相邻状态。} 规律是结果尾指针每次接较小头节点，某条链耗尽后直接接另一条剩余部分；虚拟头统一处理首节点。把这个特例继续拆开看：每个指针都要有角色，上一段尾、当前段头、当前节点和后继入口不能混。只要一次改 next 之前没有保存后继，就可能丢掉整段未处理链。

\textbf{状态复用。} 用 dummy 作为结果虚拟头，tail 指向结果尾。每次比较 l1 和 l2，接上较小节点并推进对应链表。手算时只改被新输入影响的那一小部分，而不是回头重算完整候选。

\textbf{指针和字段。} 状态字段要能说清：l1/l2 是两条未处理链的当前头，tail 是已合并链表尾部，dummy.next 是最终结果入口。手算时按这个协议跟踪：1 和 2 比较接 1；3 和 2 比较接 2；3 和 4 比较接 3；最后 l1 空，直接接剩余 4。

\textbf{代码落点。} 关键是 tail->next=chosen 后 tail=tail->next，再推进 chosen 来源链。某条链为空后 tail->next 接另一条剩余链。关键代码行必须能逐句翻译成“旧状态怎样变成新状态”，否则就只是模板。

\textbf{正确性。} 两条链都已排序，两个当前头中较小者不可能被另一条链后面的节点超过，所以它就是全局下一个节点。

\textbf{边界实验。} 先用某一条链为空、重复值、剩余整段接回、是否复用原节点做反例检查。实验时覆盖某一条链为空、重复值、剩余整段接回、是否复用原节点，每步画出前驱、当前、后继和下一段头。链表题不要只看输出值，还要检查节点是否丢失或成环。

\textbf{复杂度变化。} 排序数组法 O((m+n)log(m+n)) 且额外空间高；双指针合并 O(m+n) 时间，O(1) 额外空间。

\textbf{迁移追问。} 如果题目改成“合并 K 条链表可用分治或堆，核心仍是维护每条链的当前头”，先判断原状态是否仍然覆盖新答案集合，再决定是微调边界、改 value 语义，还是换算法。

\subsection{LeetCode 19 删除链表的倒数第 N 个结点}

\textbf{题目说明。} LeetCode 19 删除链表的倒数第 N 个结点 的题面入口要先还原成具体任务：倒数位置可以转成两个指针之间固定 n 步的距离。

\textbf{样例入口。} 拿长度 5 删除倒数第 2 个手算。快指针先走 2 步，慢指针和快指针保持固定距离；当快指针到尾后，慢指针正好停在待删节点前驱。再试删除头节点，若没有虚拟头就缺前驱。

\textbf{暴力基线。} 暴力先扫描一遍链表求长度 len，再计算要删除的是第 len-n+1 个节点；第二遍走到它的前驱后改 next。

\textbf{暴力过程。} 以 1->2->3->4->5、n=2 为例，暴力先数出长度 5，目标是正数第 4 个节点 4，再从头走到节点 3，把 3->next 改成 5。

\textbf{重复工作。} 题面模型：倒数位置可以转成两个指针之间固定 n 步的距离。暴力版的慢点要落到本题：浪费在为了定位倒数位置先完整数长度。倒数第 n 个节点和它前驱，可以通过快慢指针之间固定 n+1 步距离一次扫描定位。后面的优化只消掉这类重复工作，不能改变答案集合。

\textbf{优化条件。} 本题的关键观察是：快指针先走 n 步，再让快慢指针同步移动，慢指针停在待删节点前驱。这句话必须能翻译成可维护状态；如果题目条件改变后观察不再成立，就不能继续套原来的写法。

\textbf{相邻状态。} 规律是虚拟头统一头部删除，快慢指针固定距离定位前驱。把这个特例继续拆开看：每个指针都要有角色，上一段尾、当前段头、当前节点和后继入口不能混。只要一次改 next 之前没有保存后继，就可能丢掉整段未处理链。

\textbf{状态复用。} 使用虚拟头 dummy 指向 head。fast 先从 dummy 走 n+1 步，然后 fast 和 slow 同步走；fast 到空时，slow 正好是待删节点前驱。手算时只改被新输入影响的那一小部分，而不是回头重算完整候选。

\textbf{指针和字段。} 状态字段要能说清：dummy 负责统一删除头节点，fast 负责制造距离，slow 停在前驱，slow->next 是要删除的节点。手算时按这个协议跟踪：删除 1->2->3->4->5 的倒数第 2 个时，fast 先领先 slow 三步；同步移动到 fast 为空，slow 在 3，删除 slow->next=4。

\textbf{代码落点。} 关键代码是 dummy.next=head，fast=\&dummy，预走 n+1 步；最后 slow->next=slow->next->next。预走从 dummy 开始，才能覆盖删除头节点。关键代码行必须能逐句翻译成“旧状态怎样变成新状态”，否则就只是模板。

\textbf{正确性。} fast 和 slow 始终相隔 n+1 条边；当 fast 走到链表末尾之后，slow 后面正好还有 n 个节点，所以 slow 是目标前驱。

\textbf{边界实验。} 先用删除头节点、n 等于链表长度、只有一个节点、题目保证 n 合法做反例检查。实验时覆盖删除头节点、n 等于链表长度、只有一个节点、题目保证 n 合法，每步画出前驱、当前、后继和下一段头。链表题不要只看输出值，还要检查节点是否丢失或成环。

\textbf{复杂度变化。} 两遍扫描 O(n)；快慢指针仍是 O(n) 时间，但只走一遍主体扫描，额外空间 O(1)。

\textbf{迁移追问。} 如果题目改成“若 n 不保证合法，需要在快指针预走阶段检测长度不足”，先判断原状态是否仍然覆盖新答案集合，再决定是微调边界、改 value 语义，还是换算法。

\subsection{LeetCode 20 有效的括号}

\textbf{题目说明。} LeetCode 20 有效的括号 的题面入口要先还原成具体任务：括号匹配要求最近打开的左括号先被关闭。

\textbf{样例入口。} 拿 ([)] 手算：只计数左右括号会误判，因为闭合顺序错了。栈顶必须保存最近的未匹配左括号；遇到 ] 时栈顶若不是 [，立即失败。

\textbf{暴力基线。} 暴力反复删除字符串中的 ()、[]、\{\}，直到不能再删除；若最后为空则合法。

\textbf{暴力过程。} 以 ([)] 为例，左右括号数量都对，但没有任何相邻合法括号对可删除，说明只数数量会误判顺序。

\textbf{重复工作。} 题面模型：括号匹配要求最近打开的左括号先被关闭。暴力版的慢点要落到本题：浪费在反复扫描字符串找可删除对。括号匹配只关心最近一个尚未关闭的左括号，它天然符合后进先出。后面的优化只消掉这类重复工作，不能改变答案集合。

\textbf{优化条件。} 本题的关键观察是：栈保存尚未匹配的左括号，遇到右括号必须和栈顶类型一致。这句话必须能翻译成可维护状态；如果题目条件改变后观察不再成立，就不能继续套原来的写法。

\textbf{相邻状态。} 规律是括号匹配不是数量问题，而是最近未闭合结构的后进先出问题。把这个特例继续拆开看：栈里的对象都是答案尚未确定的候选；一旦当前元素触发弹栈，被弹出的候选必须已经同时拿到了左边界和右边界，未来元素不再有资格改变它的答案。

\textbf{状态复用。} 栈保存未匹配的左括号。遇到左括号入栈；遇到右括号时栈不能为空，且栈顶必须是对应类型，匹配后弹出。手算时只改被新输入影响的那一小部分，而不是回头重算完整候选。

\textbf{指针和字段。} 状态字段要能说清：stack.top 是最近未匹配左括号，ch 是当前字符，pairs 记录右括号需要的左括号类型。手算时按这个协议跟踪：处理 ([)]：入栈 (，入栈 [；遇到 ) 时栈顶是 [，不是 (，立即失败。

\textbf{代码落点。} 关键代码是遇到右括号先检查 empty，再比较 top。扫描结束后还必须检查栈为空，防止 (( 这类前缀合法但未闭合的输入。关键代码行必须能逐句翻译成“旧状态怎样变成新状态”，否则就只是模板。

\textbf{正确性。} 合法括号串的任意右括号都必须关闭最近的未闭合左括号；栈顶正好表示这个对象，所以每次局部匹配都是必要条件。

\textbf{边界实验。} 先用空串、以右括号开头、交叉匹配、结束后栈非空做反例检查。实验时准备空串、以右括号开头、交叉匹配、结束后栈非空，打印每次入栈、弹栈和结算对象。重点观察相等元素、空栈、哨兵和最后一次结算。

\textbf{复杂度变化。} 反复删除最坏 O(\ensuremath{n^{2}})，栈扫描 O(n) 时间，O(n) 空间。

\textbf{迁移追问。} 如果题目改成“若加入通配符星号，需要额外维护可行区间或双栈”，先判断原状态是否仍然覆盖新答案集合，再决定是微调边界、改 value 语义，还是换算法。

\subsection{LeetCode 739 每日温度}

\textbf{题目说明。} LeetCode 739 每日温度 的题面入口要先还原成具体任务：每一天等待右侧第一个更高温度。

\textbf{样例入口。} 拿 [73, 74, 75, 71, 69, 72, 76, 73] 手算，73 遇到 74 时等待 1 天；71 要等到 72 才解决。

\textbf{暴力基线。} 慢版本让每个位置向左或向右线性寻找边界，或者让每个结构反复等待匹配。它能算出答案，但很多尚未完成的候选会被未来同一个事件一起解决。放到本题就是：先按“每一天等待右侧第一个更高温度”完整试慢版本；样例里已经能看到“规律是单调递减栈保存尚未等到更高温的下标，遇到更高温时连续弹栈结算天数”，所以优化前必须先能说清慢版本枚举了哪些候选。

\textbf{暴力过程。} 拿 [73, 74, 75, 71, 69, 72, 76, 73] 手算，73 遇到 74 时等待 1 天；71 要等到 72 才解决。

\textbf{重复工作。} 题面模型：每一天等待右侧第一个更高温度。暴力版的慢点要落到本题：慢版本会围绕“每一天等待右侧第一个更高温度”反复寻找最近边界或最近未完成对象。样例规律是“规律是单调递减栈保存尚未等到更高温的下标，遇到更高温时连续弹栈结算天数”，说明旧候选一旦遇到当前输入就能被批量结算；继续为每个位置单独向左或向右扫描就是浪费。后面的优化只消掉这类重复工作，不能改变答案集合。

\textbf{优化条件。} 本题的关键观察是：单调递减栈保存尚未找到答案的下标。这句话必须能翻译成可维护状态；如果题目条件改变后观察不再成立，就不能继续套原来的写法。

\textbf{相邻状态。} 规律是单调递减栈保存尚未等到更高温的下标，遇到更高温时连续弹栈结算天数。把这个特例继续拆开看：栈里的对象都是答案尚未确定的候选；一旦当前元素触发弹栈，被弹出的候选必须已经同时拿到了左边界和右边界，未来元素不再有资格改变它的答案。

\textbf{状态复用。} 把反复寻找最近边界改成维护未完成候选；新元素只和栈顶比较，连续弹出一批已经被当前元素解决的对象。本题落点是：规律是单调递减栈保存尚未等到更高温的下标，遇到更高温时连续弹栈结算天数。手算时只改被新输入影响的那一小部分，而不是回头重算完整候选。

\textbf{指针和字段。} 状态字段要能说清：栈内对象的业务含义、当前输入、弹出时结算的对象、左边界和右边界；栈中存下标还是值要明确。手算时按这个协议跟踪：拿 [73, 74, 75, 71, 69, 72, 76, 73] 手算，73 遇到 74 时等待 1 天；71 要等到 72 才解决。然后按协议复查：手算时记录栈内元素的业务含义，不只记录数值。入栈表示答案未定，弹栈表示当前事件已经让它的答案定下来。

\textbf{代码落点。} 弹栈前确认栈非空，弹出后再读取新的左边界；相等元素和哨兵高度要有统一策略。关键代码行必须能逐句翻译成“旧状态怎样变成新状态”，否则就只是模板。

\textbf{正确性。} 证明从弹出时机开始：被弹出的对象在当前时刻答案已经确定，未来元素无法再改变它。

\textbf{边界实验。} 先用没有更高温、相等温度、递增、递减做反例检查。实验时准备没有更高温、相等温度、递增、递减，打印每次入栈、弹栈和结算对象。重点观察相等元素、空栈、哨兵和最后一次结算。

\textbf{复杂度变化。} 暴力为每个位置寻找边界常见为平方级；单调栈让每个元素最多入栈一次、出栈一次，因此主体扫描为线性时间。

\textbf{迁移追问。} 如果题目改成“下一个更大元素与它同型，只是输出值或距离不同”，先判断原状态是否仍然覆盖新答案集合，再决定是微调边界、改 value 语义，还是换算法。

\subsection{LeetCode 994 腐烂的橘子}

\textbf{题目说明。} LeetCode 994 腐烂的橘子 的题面入口要先还原成具体任务：所有初始腐烂橘子同时作为 BFS 源点。

\textbf{样例入口。} 拿 [[2, 1, 1], [1, 1, 0], [0, 1, 1]] 手算，所有腐烂橘子同时向四周扩散，第一分钟会感染相邻新鲜橘子。

\textbf{暴力基线。} 暴力按分钟模拟整张网格：每一分钟扫描所有格子，找到会被感染的新鲜橘子并统一变腐烂。

\textbf{暴力过程。} 在 [[2, 1, 1], [1, 1, 0], [0, 1, 1]] 中，初始腐烂橘子同时向四周扩散；第一分钟感染它上下左右的新鲜橘子。

\textbf{重复工作。} 题面模型：所有初始腐烂橘子同时作为 BFS 源点。暴力版的慢点要落到本题：浪费在每分钟都扫描全图。扩散只会从上一分钟刚腐烂的边界继续向外，队列正好保存这个边界。后面的优化只消掉这类重复工作，不能改变答案集合。

\textbf{优化条件。} 本题的关键观察是：按层扩散，分钟数等于最后一个新鲜橘子被首次访问的层数。这句话必须能翻译成可维护状态；如果题目条件改变后观察不再成立，就不能继续套原来的写法。

\textbf{相邻状态。} 规律是多源 BFS，初始所有腐烂橘子同时入队，层数就是分钟数。把这个特例继续拆开看：状态节点不一定等于原始位置，还可能包含钥匙、剩余资源、时间或访问集合。visited 只能合并未来能力完全相同的状态，否则会把合法路径剪掉。

\textbf{状态复用。} 多源 BFS。初始把所有腐烂橘子入队，并统计 fresh 数量；每弹出一层代表一分钟，感染相邻新鲜橘子并 fresh--。手算时只改被新输入影响的那一小部分，而不是回头重算完整候选。

\textbf{指针和字段。} 状态字段要能说清：queue 保存当前扩散边界，fresh 是剩余新鲜橘子数，minutes 是已经完成的层数，directions 是四方向。手算时按这个协议跟踪：第一层处理初始腐烂源，感染邻居入队；第二层只处理第一分钟新腐烂的橘子，层数自然对应分钟。

\textbf{代码落点。} 关键是所有初始腐烂橘子一次性入队。感染时立即把 grid[nr][nc] 改成 2，避免同一分钟被重复入队。关键代码行必须能逐句翻译成“旧状态怎样变成新状态”，否则就只是模板。

\textbf{正确性。} BFS 按层扩展，某个新鲜橘子第一次被访问的层数就是它能被腐烂的最早分钟；多源初始化表示所有源同时开始。

\textbf{边界实验。} 先用没有新鲜橘子、没有腐烂源、隔离新鲜橘子、空格阻断做反例检查。实验时覆盖没有新鲜橘子、没有腐烂源、隔离新鲜橘子、空格阻断，并打印入队时的完整状态。若同一位置但资源不同的状态被合并，很多图题会出现隐蔽错误。

\textbf{复杂度变化。} 逐分钟全图扫描最坏 O(mn*分钟数)，BFS 每格最多入队一次，O(mn) 时间、O(mn) 空间。

\textbf{迁移追问。} 如果题目改成“多源 BFS 也适用于最近零距离、地图扩散等题”，先判断原状态是否仍然覆盖新答案集合，再决定是微调边界、改 value 语义，还是换算法。

\subsection{LeetCode 102 二叉树层序遍历}

\textbf{题目说明。} LeetCode 102 二叉树层序遍历 的题面入口要先还原成具体任务：队列在每轮开始时保存当前层所有节点。

\textbf{样例入口。} 拿根 3、左 9、右 20 手算，队列第一轮只有 3，第二轮才是 9 和 20。若不固定本层 size，新入队的下一层节点会混进当前层。

\textbf{暴力基线。} 慢版本可以先给每个节点计算深度，再按深度分组；或者每层都从根重新扫描一遍收集该层节点。

\textbf{暴力过程。} 以根 3、左右孩子 9 和 20 为例，第一层只有 3，第二层是 9、20。若队列不固定当前层长度，20 的孩子可能混进同一层。

\textbf{重复工作。} 题面模型：队列在每轮开始时保存当前层所有节点。暴力版的慢点要落到本题：浪费在重复从根找层，或者层边界不清导致混层。队列天然按访问顺序保存下一批节点，只要固定本轮队列大小就是当前层。后面的优化只消掉这类重复工作，不能改变答案集合。

\textbf{优化条件。} 本题的关键观察是：记录 level size 固定层边界，避免下一层节点混入当前层。这句话必须能翻译成可维护状态；如果题目条件改变后观察不再成立，就不能继续套原来的写法。

\textbf{相邻状态。} 规律是每轮开始记录队列长度，这个长度就是当前层边界。把这个特例继续拆开看：节点向父亲返回的量和全局答案通常不是同一个量。先写叶子或空节点的返回值，再写父节点如何合并左右孩子，递归式才有边界基础。

\textbf{状态复用。} BFS 队列保存待处理节点。每轮先读取 level\_size，只弹出这 level\_size 个节点并收集值，同时把它们的孩子加入队尾。手算时只改被新输入影响的那一小部分，而不是回头重算完整候选。

\textbf{指针和字段。} 状态字段要能说清：queue 是待访问节点队列，level\_size 是当前层节点数，level\_values 是本层输出。手算时按这个协议跟踪：队列初始 [3]，level\_size=1，弹出 3 后压入 9 和 20；下一轮 level\_size=2，只处理 9 和 20。

\textbf{代码落点。} 关键是 for 循环次数使用进入本层前的 queue.size()，不要在循环条件里直接用动态变化的 queue.size()。关键代码行必须能逐句翻译成“旧状态怎样变成新状态”，否则就只是模板。

\textbf{正确性。} BFS 入队顺序保证父层节点先于子层节点；固定 level\_size 后，本轮处理的恰好是同一深度的全部节点。

\textbf{边界实验。} 先用空树、只有根、最后一层不满、左右顺序做反例检查。实验时覆盖空树、只有根、最后一层不满、左右顺序，尤其关注空节点、单节点、链式树和全负值。递归版本和显式栈版本可以互相对拍。

\textbf{复杂度变化。} 重复扫描每层最坏 O(nh)，BFS 每节点入队出队一次 O(n)，队列空间 O(width)。

\textbf{迁移追问。} 如果题目改成“锯齿形层序只是在每层输出顺序上增加方向控制”，先判断原状态是否仍然覆盖新答案集合，再决定是微调边界、改 value 语义，还是换算法。

\subsection{LeetCode 239 滑动窗口最大值}

\textbf{题目说明。} LeetCode 239 滑动窗口最大值 的题面入口要先还原成具体任务：每个窗口最大值来自仍未过期且没有被新元素支配的候选。

\textbf{样例入口。} 拿 [1, 3, -1]、窗口 3 手算，队尾的 1 在 3 进入后永远不可能成为最大值，因为 3 更新且更靠右。于是 1 可以弹出。

\textbf{暴力基线。} 暴力枚举每个长度为 k 的窗口，扫描窗口内所有元素求最大值。

\textbf{暴力过程。} 以 [1, 3, -1, -3, 5]、k=3 为例，第一个窗口 [1, 3, -1] 最大值为 3；下一个窗口 [3, -1, -3] 仍要重新扫描会重复看 3 和 -1。

\textbf{重复工作。} 题面模型：每个窗口最大值来自仍未过期且没有被新元素支配的候选。暴力版的慢点要落到本题：浪费在相邻窗口重复求最大值。若新来的数更大且位置更靠右，队尾所有不大于它的旧数未来都不可能再当最大值。后面的优化只消掉这类重复工作，不能改变答案集合。

\textbf{优化条件。} 本题的关键观察是：单调队列保存下标，队尾小于等于新值的候选被弹出。这句话必须能翻译成可维护状态；如果题目条件改变后观察不再成立，就不能继续套原来的写法。

\textbf{相邻状态。} 规律是单调队列保存可能成为当前或未来窗口最大值的下标，过期下标从队首清理。把这个特例继续拆开看：栈里的对象都是答案尚未确定的候选；一旦当前元素触发弹栈，被弹出的候选必须已经同时拿到了左边界和右边界，未来元素不再有资格改变它的答案。

\textbf{状态复用。} 单调递减队列保存下标。入队前从队尾弹出 nums[back]<=nums[i] 的下标；队首若过期 i-k 则弹出；窗口形成后队首就是最大值。手算时只改被新输入影响的那一小部分，而不是回头重算完整候选。

\textbf{指针和字段。} 状态字段要能说清：deque.front 是当前最大值下标，i 是右端，i-k+1 是当前窗口左端。手算时按这个协议跟踪：1 入队后，3 进入时弹出 1，因为 3 更大且更晚；窗口形成后队首 3 给出最大值。

\textbf{代码落点。} 关键是队列保存下标而不是值，这样才能判断过期。顺序通常是清队尾、入队、清队首过期、读答案。关键代码行必须能逐句翻译成“旧状态怎样变成新状态”，否则就只是模板。

\textbf{正确性。} 被队尾新元素弹出的旧元素既更小又更早，在未来任何包含旧元素的窗口里也会包含新元素，因此不会成为最大值。

\textbf{边界实验。} 先用k 为一、k 等于数组长度、重复最大值、队首过期做反例检查。实验时准备k 为一、k 等于数组长度、重复最大值、队首过期，打印每次入栈、弹栈和结算对象。重点观察相等元素、空栈、哨兵和最后一次结算。

\textbf{复杂度变化。} 暴力 O(nk)，单调队列每个下标入队出队一次，O(n) 时间、O(k) 空间。

\textbf{迁移追问。} 如果题目改成“受限子序列和使用同样队列维护最近 k 个 DP 最大值”，先判断原状态是否仍然覆盖新答案集合，再决定是微调边界、改 value 语义，还是换算法。

\subsection{LeetCode 53 最大子数组和}

\textbf{题目说明。} LeetCode 53 最大子数组和给定整数数组 nums，要求找一个非空连续子数组，使元素和最大，并返回这个最大和。

\textbf{样例入口。} 样例 nums=[-2, 1, -3, 4, -1, 2, 1, -5, 4]，输出 6；连续段 [4, -1, 2, 1] 的和为 6。

\textbf{暴力基线。} 暴力枚举所有连续子数组，累计每段和并更新最大值；优化一点可以固定左端、右端扩展时维护当前和。

\textbf{暴力过程。} 以 [-2, 1, -3, 4, -1, 2, 1] 为例，暴力会重复计算很多以 4 开头或以 2 结尾的子段。真正关键是当前数要不要接上前一个后缀。

\textbf{重复工作。} 题面模型：给定整数数组，返回一个非空连续子数组能够取得的最大元素和。暴力版的慢点要落到本题：浪费在重复求每个后缀和。对于以 i 结尾的最优连续段，历史信息只需要一个值：以 i-1 结尾的最大后缀和。后面的优化只消掉这类重复工作，不能改变答案集合。

\textbf{优化条件。} 本题的关键观察是：用 [-2, 3] 和 [2, 3] 对比，推出当前位置只在单独开始与接上旧后缀之间取较大值。这句话必须能翻译成可维护状态；如果题目条件改变后观察不再成立，就不能继续套原来的写法。

\textbf{相邻状态。} 规律是当前位置只需要比较 nums[i] 与 bestEndingHere + nums[i]，全局答案再和 bestEndingHere 比较；全负数组要求答案初始化为第一个数，不能初始化为 0。把这个特例继续拆开看：先问暴力搜索里哪些不同路径到达同一个后续问题，再给这个后续问题命名为状态。状态必须包含未来决策需要的全部信息，转移只从更小且已经算清的状态来。

\textbf{状态复用。} bestEndingHere 表示以当前位置结尾的最大子数组和；answer 表示全局最大值。当前位置要么单独开始，要么接上旧后缀。手算时只改被新输入影响的那一小部分，而不是回头重算完整候选。

\textbf{指针和字段。} 状态字段要能说清：current=bestEndingHere，answer 是全局最优；nums[i] 与 current+nums[i] 比较决定是否丢弃旧后缀。手算时按这个协议跟踪：遇到 4 之前，旧后缀为负，接上反而变差，所以从 4 重新开始；之后 -1、2、1 继续接上得到 6。

\textbf{代码落点。} 关键代码是 current=max(nums[i], current+nums[i]); answer=max(answer, current)。answer 必须初始化为 nums[0]，否则全负数组会错。关键代码行必须能逐句翻译成“旧状态怎样变成新状态”，否则就只是模板。

\textbf{正确性。} 任意以 i 结尾的连续段只有两类：只含 nums[i]，或由某个以 i-1 结尾的连续段加上 nums[i]。取前一类最优即可覆盖全部候选。

\textbf{边界实验。} 先用全负数组、单元素、和溢出、答案不能初始化为零做反例检查。实验时覆盖全负数组、单元素、和溢出、答案不能初始化为零，先打印完整 DP 表，再比较空间压缩版本。若压缩版不同，优先检查遍历顺序和旧值保存。

\textbf{复杂度变化。} 暴力 O(\ensuremath{n^{2}}) 或 O(\ensuremath{n^{3}})，Kadane 算法 O(n) 时间、O(1) 空间。

\textbf{迁移追问。} 如果题目改成“若要求返回区间下标，同步记录重新开始位置和最优左右端”，先判断原状态是否仍然覆盖新答案集合，再决定是微调边界、改 value 语义，还是换算法。

\subsection{LeetCode 76 最小覆盖子串}

\textbf{题目说明。} LeetCode 76 最小覆盖子串给定字符串 s 和 t，要求在 s 中找最短连续子串，使它包含 t 中每个字符及其次数；不存在则返回空串。

\textbf{样例入口。} 样例 s=ADOBECODEBANC，t=ABC，输出 BANC；它是覆盖 A、B、C 的最短连续子串。

\textbf{暴力基线。} 暴力枚举 s 的每个子串，重新统计字符频次，检查是否覆盖 t 的每个字符及其次数。

\textbf{暴力过程。} 以 s=ADOBECODEBANC、t=ABC 为例，暴力会检查 A、AD、ADO...直到 ADOBEC 才第一次覆盖；下一轮从 D 开始又重新统计 DOBEC。

\textbf{重复工作。} 题面模型：给定字符串 s 和 t，在 s 中找一个最短连续子串，使它覆盖 t 的每个字符及其次数。暴力版的慢点要落到本题：浪费在相邻子串重复统计字符。右端扩展只新增一个字符，左端收缩只移走一个字符。后面的优化只消掉这类重复工作，不能改变答案集合。

\textbf{优化条件。} 本题的关键观察是：窗口内字符计数必须覆盖目标计数，满足后尽量收缩左端；用 formed 记录满足需求的字符种类，避免每次全量比较。这句话必须能翻译成可维护状态；如果题目条件改变后观察不再成立，就不能继续套原来的写法。

\textbf{相邻状态。} 规律是右端负责获得可行，左端负责在可行时尽量缩短；计数表要区分已有字符和需求字符。把这个特例继续拆开看：右端进入只带来一个新元素，左端离开只删掉一个旧元素；只有当旧左端被移走后不可能再产生更优答案时，窗口才可以单向推进。若这个单调淘汰条件不存在，就要换成前缀和、队列或其他状态。

\textbf{状态复用。} need 保存目标频次，window 保存当前窗口频次，formed 或 required\_count 记录当前满足需求的字符种类；满足后尽量收缩 left。手算时只改被新输入影响的那一小部分，而不是回头重算完整候选。

\textbf{指针和字段。} 状态字段要能说清：right 负责扩展获得覆盖，left 负责在覆盖时缩短窗口，best\_left/best\_len 记录最短答案。手算时按这个协议跟踪：ADOBEC 第一次覆盖 ABC，记录长度 6；左端移走 A 后不覆盖，停止收缩。继续扩展到 BANC 时再次覆盖，收缩得到长度 4。

\textbf{代码落点。} 关键顺序是加入 s[right] 后更新 formed；while formed==required 时先更新答案，再移走 s[left] 并可能降低 formed。关键代码行必须能逐句翻译成“旧状态怎样变成新状态”，否则就只是模板。

\textbf{正确性。} 每个 right 下，while 循环会把 left 推到刚好不能再缩的位置，因此不会错过以该 right 结尾的最短覆盖窗口。

\textbf{边界实验。} 先用目标有重复字符、源串缺字符、最小窗口在开头或结尾、大小写做反例检查。实验时重点跑目标有重复字符、源串缺字符、最小窗口在开头或结尾、大小写，并打印左右端、窗口计数和答案。若左端发生回退，或者窗口失去合法性后仍更新答案，就能很快定位错误。

\textbf{复杂度变化。} 暴力至少 O(\ensuremath{n^{2}}*字符集)，滑动窗口每个字符进出一次，时间 O(n+字符集)，空间 O(字符集)。

\textbf{迁移追问。} 如果题目改成“若字符集固定，数组计数更快；若是 Unicode，哈希表更通用”，先判断原状态是否仍然覆盖新答案集合，再决定是微调边界、改 value 语义，还是换算法。

\begin{principlebox}{本节复盘方式}
不要只看最终算法名。每道题复盘时先遮住优化版，口述暴力枚举对象；再指出相邻窗口、历史前缀、候选区间、搜索状态或子问题里哪一部分被重复处理；最后用一个边界样例验证状态更新顺序。能讲完这三步，才算真正掌握本章案例。
\end{principlebox}
